\documentclass[10pt,journal,onecolumn]{IEEEtran}

% ---------- Packages ----------
\usepackage[T1]{fontenc}
\usepackage{newtxtext}
\usepackage{amsmath}
\usepackage{amssymb}
\usepackage{newtxmath}
\usepackage{booktabs}
\usepackage{array}
\usepackage{tabularx}
\usepackage{listings}
\usepackage{xcolor}
\usepackage[hidelinks]{hyperref}
\usepackage{cleveref}

% ---------- Paragraph breaking ----------
\emergencystretch=3em
\tolerance=2000

% ---------- Code style ----------
\definecolor{codegray}{gray}{0.94}
\definecolor{codebar}{gray}{0.65}
\lstdefinestyle{sec}{
  basicstyle=\ttfamily\scriptsize,
  backgroundcolor=\color{codegray},
  breaklines=true,
  showstringspaces=false,
  frame=single,
  rulecolor=\color{codebar},
  framerule=0.3pt,
  columns=fullflexible,
  keepspaces=true,
  aboveskip=4pt,
  belowskip=4pt
}
\lstset{style=sec}

\newcommand{\std}[1]{\textit{#1}}

% ---------- Title / Authors ----------
\title{Compliance and Secure Coding Practices Report:\\[0.3em]
NeuroShell: An Autonomous Natural-Language-Driven Orchestration
Framework for Offensive Security Operations}
\author{\IEEEauthorblockN{NeuroShell Core Team --- Group R26-CS-018}
\IEEEauthorblockA{Department of Computer Systems Engineering\\
MSc Cybersecurity Research Project\\
Supervisor: Prof.\ Harinda Fernando}}

\begin{document}

\maketitle

\begin{abstract}
This report documents the compliance and secure coding practices applied
throughout the development of NeuroShell, a four-component framework
(Intent Recognition Engine, Dynamic Planner/RAG, Adaptive Execution and
Error Recovery Engine, and AI-driven Vulnerability Analysis Engine) for
authorised offensive-security operations. It states the compliance
standards adopted, the secure coding standards followed, how both were
embedded in the Secure Development Lifecycle (SDLC), the tools used, the
code review and audit process, the ethical safeguards, and the principal
challenges and their resolutions. All statements are grounded in the
project's verification artefacts (506 automated tests, 76\% statement
coverage, CI security scanning, and the append-only audit log).
\end{abstract}

\begin{IEEEkeywords}
Compliance, secure coding, GDPR, ISO/IEC 27001, OWASP, CERT, SSDF,
secure SDLC, prompt injection, RBAC, audit logging.
\end{IEEEkeywords}

% ============================================================
\section{Compliance Standards}
% ============================================================

\subsection{Applicable and non-applicable regimes}
NeuroShell is an authorised-assessment tool operated on private
infrastructure under a defined engagement scope. It does not process
health information, so \std{HIPAA} is not applicable. It processes
operator commands and the technical artefacts of an assessment rather
than personal data of data subjects, so the applicable regimes are data
protection law as far as operator records are concerned, and information
security management for the software itself. \cref{tab:compliance} maps
each standard to its applicability and to the mechanism that implements
it.

\begin{table}[htbp]
  \centering
  \footnotesize
  \begin{tabular}{@{}p{0.24\textwidth}p{0.20\textwidth}p{0.46\textwidth}@{}}
    \toprule
    \textbf{Standard} & \textbf{Applicability} & \textbf{How it was integrated}\\
    \midrule
    GDPR (EU 2016/679) & Partial --- operator audit records & Data minimisation and pseudonymisation: audit log stores a SHA-256 digest, never the raw command; purpose and storage limitation enforced by policy\\
    ISO/IEC 27001:2022 & Applicable --- ISMS for the software & Access control (A.9), operations security (A.12), secure development (A.14) reflected in RBAC, logging, and CI controls\\
    HIPAA & Not applicable & No protected health information is collected or processed\\
    \bottomrule
  \end{tabular}
  \caption{Compliance standards mapped to their project implementation.}\label{tab:compliance}
\end{table}

\subsection{GDPR integration}
The GDPR principles were applied to the one place where operator input is
persisted: the audit log. Raw natural-language commands are never stored;
a SHA-256 digest is retained instead (\cref{lst:hash}), which satisfies
data minimisation (Art.\ 5(1)(c)) and the security-of-processing
obligation (Art.\ 32) while keeping the record useful for
accountability. Retention is bounded by policy, and the audit store is
append-only, giving integrity without exposing input content.

\subsection{ISO/IEC 27001 integration}
ISO/IEC 27001 provided the control framework for the software itself.
Least-privilege access control is implemented by the role hierarchy
(viewer $\rightarrow$ analyst $\rightarrow$ operator $\rightarrow$ admin;
\cref{lst:rbac}), where an unknown role is treated as the most
restrictive. Operations security is evidenced by structured logging,
health checks, and an append-only audit trail, and secure development is
evidenced by the CI gates described in \cref{sec:review}.

% ============================================================
\section{Secure Coding Standards}
% ============================================================

\subsection{OWASP Secure Coding Practices and Top 10}
The OWASP Secure Coding Practices and the OWASP Top 10 (2021) guided the
input-handling and access-control code, and the OWASP Top 10 for LLM
Applications (2025) guided the language-model layer. \cref{tab:owasp}
maps the highest-risk categories to concrete controls.

\begin{table}[htbp]
  \centering
  \footnotesize
  \begin{tabular}{@{}p{0.34\textwidth}p{0.58\textwidth}@{}}
    \toprule
    \textbf{Risk category} & \textbf{Control in NeuroShell}\\
    \midrule
    LLM01 Prompt injection & Three-layer AdversarialDetector (lexical, structural, scoring) with auto-block on critical\\
    A03 Injection & RegexValidator shell-metacharacter scan and parameterised SQLite statements\\
    A01 Broken access control & RBACGuard enforces a pre/post-inference role matrix; least privilege by default\\
    A02 Cryptographic failures & SHA-256 for audit integrity; HS256 JWT with short 30-minute expiry\\
    A09 Logging failures & Structured logging plus append-only, hashed audit records\\
    \bottomrule
  \end{tabular}
  \caption{OWASP risk categories mapped to NeuroShell controls.}\label{tab:owasp}
\end{table}

\subsection{CERT and ISO/IEC 27001 coding expectations}
Following the SEI CERT secure-coding guidance, the codebase avoids
dynamic execution constructs (\texttt{eval}/\texttt{exec}), validates
every external field with explicit regular expressions and Pydantic
schemas, bounds resources through rate limiting, and handles failures
without leaking internal state. Metadata from the language model is never
trusted: every entity it produces (IP addresses, CIDR blocks, CVE
identifiers, ports, URLs) is re-verified deterministically before use.

\begin{lstlisting}[caption={Privacy-preserving, parameterised audit write (\texttt{audit\_logger.py}).},label={lst:hash}]
def hash_raw_input(raw_input: str) -> str:
    return hashlib.sha256(raw_input.encode("utf-8")).hexdigest()
...
cursor.execute("INSERT INTO audit_log (...) VALUES (?, ?, ?, ...);",
               (ts, session_id, raw_hash, intent_summary, ...))
\end{lstlisting}

\begin{lstlisting}[caption={Least-privilege role gating (\texttt{rbac\_policy.py}).},label={lst:rbac}]
INTENT_ROLE_POLICY = {
    IntentType.EXPLOITATION: ["operator", "admin"],
    IntentType.NETWORK_SCAN: ["analyst", "operator", "admin"], ...}
# Unknown roles are treated as viewer (most restrictive).
def is_role_permitted(role, intent):
    return role in INTENT_ROLE_POLICY.get(intent, [])
\end{lstlisting}

% ============================================================
\section{Integration of the Secure Development Lifecycle}
% ============================================================
Security and compliance were treated as first-class requirements rather
than a post-hoc review. \cref{tab:sdlc} summarises the emphasis on each
SDLC phase.

\begin{table}[htbp]
  \centering
  \footnotesize
  \begin{tabular}{@{}p{0.20\textwidth}p{0.72\textwidth}@{}}
    \toprule
    \textbf{Phase} & \textbf{Security and compliance emphasis}\\
    \midrule
    Requirements & Safety invariant and functional/non-functional security requirements (RBAC roles, scope containment, rate limits) specified up front\\
    Design & Defence in depth: independent deterministic validators that cannot be bypassed by the model; contract schema fixed before inference\\
    Implementation & Type-annotated modules, input normalisation, no dynamic execution, secrets kept out of source\\
    Testing & 506 deterministic tests (434 unit, 72 integration); adversarial, injection, and RBAC negative paths\\
    Deployment & Pinned dependencies, minimal container images, environment-injected secrets, health checks\\
    Maintenance & CI gates (lint, type check, test, Bandit security scan) on every pull request\\
    \bottomrule
  \end{tabular}
  \caption{Security and compliance emphasis across the SDLC phases.}\label{tab:sdlc}
\end{table}

% ============================================================
\section{Tools for Security and Compliance}
% ============================================================
A layered toolset supports both security analysis and compliance
management (\cref{tab:tools}). Static application security testing is
performed by Bandit, which scans the codebase for common Python
weaknesses and is run as a dedicated CI job. Linting, formatting, and
type checking (Ruff, Black, Mypy) keep the code consistent and catch
defects early, while Pytest with Coverage measures the verification
gate. GitHub Actions orchestrates the pipeline, and Docker with the Kong
API gateway provide reproducible, isolated deployment.

\begin{table}[htbp]
  \centering
  \footnotesize
  \begin{tabular}{@{}p{0.22\textwidth}p{0.30\textwidth}p{0.40\textwidth}@{}}
    \toprule
    \textbf{Tool} & \textbf{Purpose} & \textbf{Role in the project}\\
    \midrule
    Bandit & SAST & Security scan in CI (\texttt{-x ./tests,./infra})\\
    Ruff / Black & Linting / format & Style and defect drift control in CI\\
    Mypy & Type checking & Static type verification of services\\
    Pytest + Coverage & Testing & 506 tests passing, 76\% coverage\\
    GitHub Actions & Orchestration & Lint, typecheck, test, security-scan, build\\
    Docker / Kong & Isolation / gateway & Reproducible services and routed edge\\
    \bottomrule
  \end{tabular}
  \caption{Security and compliance tooling.}\label{tab:tools}
\end{table}

% ============================================================
\section{Code Review and Audit Process}\label{sec:review}
% ============================================================
All changes reach the protected branches only through pull requests, and
the CI pipeline enforces the review gate: the pull request cannot merge
unless the lint, type-check, test, and security-scan jobs pass, and the
container build is also exercised for reviewed changes. Reviewers check
both correctness and security-relevant behaviour, with particular
attention to the deterministic validators, the RBAC matrix, and any code
touching persisted data. Feedback is incorporated in subsequent commits
on the same branch, and additional regression tests are added when a
review uncovers a weakness. The audit dimension is satisfied at runtime
by an append-only SQLite trail: every enforcement decision is recorded
with its class, severity, and action, raw input is hashed, and unresolved
findings enter an escalation queue that an administrator resolves through
a dedicated endpoint. A Quality Assurance report, produced from the test,
coverage, and static-analysis runs, captures the outcome and feeds the
next development cycle.

% ============================================================
\section{Ethical and Compliance Considerations}
% ============================================================
NeuroShell is a dual-use offensive-security tool, so its ethical
guardrails are enforced in code rather than left to operator discretion.
Targets are constrained to the RFC\,1918 private-address space and the
configured engagement scope, so the framework cannot be aimed at systems
outside the authorised boundary. High-impact operations are gated by the
role hierarchy, and ambiguous or contradictory requests are refused or
routed to a human reviewer instead of being executed automatically. No
personal data of data subjects is collected, operator input is
pseudonymised before storage, and all usage is subject to the project's
academic-use and responsible-disclosure commitments. AI assistance in the
project is documented separately in the AI-use disclosure appendix of the
thesis.

% ============================================================
\section{Challenges and Solutions}
% ============================================================
Implementing secure coding and compliance in an LLM-driven system
presented several recurring challenges; \cref{tab:challenges} records
each and how it was addressed.

\begin{table}[htbp]
  \centering
  \footnotesize
  \begin{tabular}{@{}p{0.40\textwidth}p{0.52\textwidth}@{}}
    \toprule
    \textbf{Challenge} & \textbf{Solution adopted}\\
    \midrule
    Prompt injection and model hallucination & Deterministic validation stack and adversarial detector that override the model; never trust model output\\
    Balancing security strictness against usability & Configurable enforcement modes (research vs.\ strict) and a warn/block distinction per finding class\\
    Privacy of operator input while keeping auditability & Store SHA-256 digests only, in an append-only log, and never persist raw commands\\
    Ensuring consistent checks across a multi-service codebase & Central CI gates (lint, type, test, Bandit) and a shared, typed contract between components\\
    Verifying safety without a live model & Deterministic test fixtures that mock the inference engine, so the safety core is tested independently\\
    \bottomrule
  \end{tabular}
  \caption{Challenges encountered and the solutions adopted.}\label{tab:challenges}
\end{table}

\section*{Conclusion}
Compliance and secure coding were embedded throughout NeuroShell rather
than added at the end. GDPR principles shape the privacy-preserving audit
log, ISO/IEC 27001 controls the access and operations posture, and the
OWASP, CERT, and NIST SSDF guidance directs the input handling, LLM
hardening, and CI verification gates. The result is a framework whose
safety-critical behaviour is deterministic, testable, and auditable, and
whose residual risks are documented rather than hidden.

% ============================================================
\begin{thebibliography}{9}
\footnotesize
\bibitem{owasp-scp} OWASP Foundation, \emph{OWASP Secure Coding Practices Quick Reference Guide}, 2010.
\bibitem{owasp-top10} OWASP Foundation, \emph{OWASP Top 10:2021 --- The Ten Most Critical Web Application Security Risks}, 2021.
\bibitem{owasp-llm} OWASP Foundation, \emph{OWASP Top 10 for Large Language Model Applications}, 2025.
\bibitem{cert} CERT Division, Software Engineering Institute, Carnegie Mellon University, \emph{SEI CERT Coding Standards}, 2023.
\bibitem{iso27001} International Organization for Standardization, \emph{ISO/IEC 27001:2022 --- Information security, cybersecurity and privacy protection --- Information security management systems --- Requirements}, 2022.
\bibitem{gdpr} European Parliament and Council, \emph{Regulation (EU) 2016/679 (General Data Protection Regulation)}, 2016.
\bibitem{nist-ssdf} National Institute of Standards and Technology, \emph{NIST SP 800-218: Secure Software Development Framework (SSDF) Version 1.1}, 2022.
\bibitem{bandit} PyCQA, \emph{Bandit: A security linter for Python}. [Online]. Available: \url{https://bandit.readthedocs.io}
\end{thebibliography}

\end{document}
